% short/04_physics.tex: why the hand is detectable: the mechanism (the mirror, the swing of the forearm) and the
% evidence (Figure 6, the average curves; Figure 7, the first-half yaw), with the AUC explained where it is first used.
% Uses the gravity frame, the yaw, the roll and the 0-100 % time axis of Section 3. The two figures tell the hands
% apart by line style as well as by colour (black-and-white print).
\section{Why the hand is detectable}\label{sec:physics}

\begin{figure}[tbp]
\centering
\begin{minipage}[t]{0.49\linewidth}\centering
  \genfig{core/S2_23_mean_curves}
  \caption{The average movement to the nose with each hand (left: red, dashed; right: blue, solid): each
  participant's mean curve, averaged over the \R{describe/n_persons} participants (band: $\pm 1$ standard error).
  Time runs from the start of the movement (0\,\%) to the arrival at the nose (100\,\%). Top: tilting speed; middle:
  yaw rate (turning speed about the vertical); bottom: the yaw, with the first half of the movement shaded.}
  \label{fig:means}
\end{minipage}\hfill
\begin{minipage}[t]{0.49\linewidth}\centering
  \genfig{core/S2_24_first_half_yaw}
  \caption{The first-half yaw of all \R{describe/primary_trials} trials (\R{describe/primary_left} left-hand,
  hatched; \R{describe/primary_right} right-hand, filled): how far the phone turned about the vertical in the first
  half of the movement to the nose (positive: counter-clockwise seen from above). Dashed: the medians of the two
  hands. Trials beyond $\pm 45^\circ$ are drawn at the edge.}
  \label{fig:fhy}
\end{minipage}
\end{figure}

\paragraph{The mechanism.} Reflecting a movement left to right, as in a mirror, gives the same movement made with the
other hand. The reflection keeps the speed, the duration, every up--down and towards-the-face motion and the tilting
forward and back, but it reverses every sideways motion, every turn about the vertical (clockwise becomes
counter-clockwise) and every roll, the turn about the towards-the-face direction. Because the two arms are close to
mirror images, the hand can show mainly in such reversed quantities, and the turn about the vertical is the most
natural one. When the phone is brought to the nose, the forearm swings towards the body's midline, so, seen from
above, the phone turns counter-clockwise in the right hand and clockwise in the left. Near the nose the wrist turns the
phone partly back to face the person, so the turn in the first half of the movement should carry the clearest signal.

\paragraph{The evidence.} The average movements of the \R{describe/n_persons} participants show exactly this
(Figure~\ref{fig:means}). The phone tilts at the same speed with both hands (top panel), but it turns about the
vertical in opposite directions over most of the first half of the movement, counter-clockwise with the right hand and
clockwise with the left (middle); added up, the yaw of the two hands ends with opposite signs (bottom). To summarise
each trial by one number, we take its \emph{first-half yaw}: the yaw at the halfway point in time, that is, the yaw
rate added up from the start of the movement until halfway to the nose. Its median is
\R{describe/fhy_median_left_deg}$^\circ$ for the left hand and $+$\R{describe/fhy_median_right_deg}$^\circ$ for the
right (Figure~\ref{fig:fhy}), and its sign alone (positive: right hand), with no training at all, names the hand in
\R{describe/fhy_sign_rule_accuracy} of the trials. The two histograms overlap, however, so on the remaining trials the
sign names the wrong hand: one number is not enough.

How well a number separates the hands is measured by the AUC, the area under the ROC curve (which traces the share of
right-hand trials called right against the share of left-hand trials called right as the decision threshold
moves). Its meaning is simple: the chance that a randomly chosen right-hand trial has a larger value than a
randomly chosen left-hand trial; it is 0.5 for a useless number and 1
for a perfect one. For the first-half yaw the AUC is \R{describe/fhy_auc_pooled}. Within each participant it is above
0.5 for all \R{describe/fhy_persons_auc_above_half} participants (median \R{describe/fhy_person_auc_median}), and it is
about the same on Android phones (\R{describe/fhy_auc_android}) and on iPhones (\R{describe/fhy_auc_ios}). The yaw at
the end of the movement separates the hands less well (AUC \R{describe/net_yaw_auc_pooled}), because the wrist turns the
phone back near the nose. Numbers that a mirror leaves unchanged carry almost no information about the hand: of the
\R{m1/cat_n_features} single-number summaries of the movement that we examined, the \R{m1/cat_n_even} that the mirror
leaves unchanged, such as the peak rotation speed (AUC \R{m1/cat_peak_gyro_auc_2}) and the duration
(\R{m1/cat_duration_s_auc_2}), have AUCs between \R{m1/cat_even_auc_min} and \R{m1/cat_even_auc_max}, while the
\R{m1/cat_n_odd} that the mirror reverses reach up to \R{m1/cat_odd_auc_max} (for example the yaw one second after the start
of the movement), a little above the first-half yaw.
